\section{Spectral entropy and sharp occupation bounds}\label{sec:entropy61}
The pointwise mixing block in Lemma~\ref{lem:orbitgap57} costs a
logarithm at each selected cut. We now replace that block argument by an
entropy estimate along the selected cuts themselves. The auxiliary
entropy is attached to a smoothed law of \emph{centroid directions}; it
is not the Shannon entropy of the hidden register. The normalization by
centroid mass will be essential.

Let $E$ have real dimension $m\ge2$, let $\Sigma=\{u\in E:\norm u=1\}$,
and let $\sigma$ be normalized round measure on $\Sigma$. Write
$d_\Sigma$ for its geodesic distance and $W_2$ for the corresponding
quadratic Wasserstein distance. All logarithms in this section are
natural. Let $G\subseteq O(E)$ be compact, with normalized Haar measure
$m_G$, and let $\lambda$ be a probability on the command letters. On
$L^2(\Sigma,\sigma)$ define
\[
 Pf(x)=\sum_a\lambda_a f(R(u_a)^{-1}x),\qquad
 \Pi f(x)=\int_G f(R(g)^{-1}x)\,dm_G(g).
\]
Here $R$ denotes the given action; writing $G\subseteq O(E)$ identifies
its image when convenient. The projection $\Pi$ need not have
one-dimensional range. Our spectral hypothesis is
\begin{equation}\label{eq:actiongap61}
 \norm{Pf-\Pi f}_2\le\kappa\norm{f-\Pi f}_2,
                         \qquad 0\le\kappa<1.
\end{equation}
This is a bound on the whole function space, not merely on the
finite-dimensional representation $E$. The group bound
\eqref{eq:groupgap57} implies \eqref{eq:actiongap61}, with the same
$\kappa$: decompose the unitary representation on $L^2(\Sigma)$ into
irreducibles; each nontrivial irreducible occurs in the regular
representation, where the assumed bound holds (use the adjoint
operator if needed for the inverse convention). The trivial summands
are exactly the range of $\Pi$ \cite{Folland}.

\begin{remark}[The full spectral hypothesis]\label{rem:fullgap61}
A contraction on $E$ alone does not imply \eqref{eq:actiongap61}.
For finitely many circle rotations, simultaneous rational approximation
provides arbitrarily high characters on which the averaging eigenvalue
is arbitrarily close to one, although its first harmonic may contract.
Such an alphabet does not satisfy the full spectral hypothesis. This
is why the present theorem does not replace the different Diophantine
width laws by a spurious spherical exponent.
\end{remark}

\subsection{A compactly supported entropy potential}
For $0<h\le1$ set
\begin{equation}\label{eq:kernel61}
 K_h(u,x)=Z_h^{-1}\left(1-\frac{\norm{x-u}^2}{h^2}\right)_+^2,
 \qquad Z_h=\int_\Sigma\left(1-\frac{\norm{x-u}^2}{h^2}\right)_+^2d\sigma(x).
\end{equation}
The normalizing constant is independent of $u$. If $\nu$ is a
probability on $\Sigma$, write $K_h\nu(x)=\int K_h(u,x)d\nu(u)$.
For a density $f$ put
\[
 \mathcal H(f)=\int f\log f\,d\sigma,\qquad
 \mathcal I(f)=\int_{\{f>0\}}\frac{\norm{\nabla f}^2}{f}\,d\sigma,
 \qquad 0\log0=0.
\]

\begin{lemma}[Entropy continuity for spherical mixtures]\label{lem:smoothentropy61}
There are constants $D_m\ge1$ and $B_m<\infty$, depending only on $m$,
such that every finitely supported probability $\nu$ satisfies
\begin{equation}\label{eq:kernelbounds61}
 0\le\mathcal H(K_h\nu)\le L_h:=\log(D_mh^{-(m-1)}),\qquad
 \mathcal I(K_h\nu)\le B_mh^{-2}.
\end{equation}
For any two finitely supported probabilities $\nu,\eta$,
\begin{equation}\label{eq:enttransport61}
 |\mathcal H(K_h\nu)-\mathcal H(K_h\eta)|
                 \le \frac{\sqrt{B_m}}h W_2(\nu,\eta).
\end{equation}
For $m=3$ one can take $D_3=12$ and $B_3=24$.
\end{lemma}
\begin{proof}
Put $s=\norm{x-u}^2$. Polar coordinates give a density proportional to
$s^{(m-3)/2}(1-s/4)^{(m-3)/2}$ on $0<s<4$. On $0<s\le1$ its second
factor is bounded above and below by positive dimension-dependent
constants. Substitution $s=h^2v$ therefore gives
$Z_h\asymp_m h^{m-1}$. Also $\norm{\nabla s}^2=4s-s^2$, so
\[
 \mathcal I(K_h(u,\cdot))=
 \frac4{Z_hh^4}\int_{s<h^2}(4s-s^2)\,d\sigma
                            \le B_mh^{-2}.
\]
Increasing $D_m$ gives $K_h\nu\le D_mh^{-(m-1)}$.
Jensen's inequality gives the lower entropy bound, and the pointwise
upper bound gives the upper one. Cauchy--Schwarz in a finite mixture
shows that its Fisher information is at most the mixture of the
individual Fisher informations.

Choose an optimal finite coupling $(\pi_{ij})$ of $\nu$ and $\eta$.
Join its paired centers by shortest constant-speed spherical geodesics.
Each path is $u_{ij}(t)=\exp(tA_{ij})u_i$, where $A_{ij}$ is
skew-symmetric, acts on a two-dimensional plane, and has operator norm
$d_\Sigma(u_i,v_j)$. An antipodal pair is treated by choosing any such
plane. Define
\[
 f_t=\sum_{i,j}\pi_{ij}K_h(u_{ij}(t),\cdot),\qquad
 J_t(x)=\sum_{i,j}\pi_{ij}K_h(u_{ij}(t),x)A_{ij}x.
\]
Rotation invariance gives $\partial_t f_t+\operatorname{div}J_t=0$.
Pointwise Cauchy--Schwarz and integration yield
\[
 \int_{\{f_t>0\}}\frac{\norm{J_t}^2}{f_t}\,d\sigma
 \le\sum_{i,j}\pi_{ij}d_\Sigma(u_i,v_j)^2=W_2(\nu,\eta)^2.
\]
The preceding Fisher bound holds at every $t$. To justify entropy
differentiation even where $f_t$ vanishes, replace $f_t$ by
$(f_t+\delta)/(1+\delta)$ and $J_t$ by $J_t/(1+\delta)$.
The kernel is $C^1$, the sphere is compact, and integration by parts
followed by Cauchy--Schwarz gives
\[
 \left|\frac d{dt}\mathcal H\left(\frac{f_t+\delta}{1+\delta}\right)\right|
 \le\frac1{1+\delta}\sqrt{\mathcal I(f_t)}
       \left(\int\frac{\norm{J_t}^2}{f_t}\,d\sigma\right)^{1/2}
 \le\frac{\sqrt{B_m}}h W_2(\nu,\eta).
\]
Here quotients at common zeros are zero; the displayed regularized
identity can equivalently be obtained by smoothing the $C^1$ kernel.
Integrate in $t$ and let $\delta\downarrow0$. Uniform boundedness and
continuity of $x\log x$ at zero justify the limit and prove
\eqref{eq:enttransport61}.

On $\Sigma=S^2$, the variable $s$ has constant density $1/4$.
Direct integration gives
\[
 Z_h=h^2/12,\qquad
 \mathcal I(K_h(u,\cdot))=24/h^2-4.
\]
These imply the asserted explicit constants.
\end{proof}
The relation between entropy continuity and quadratic transport is
classical; compare Polyanskiy--Wu \cite[Proposition~1]{PW61} for regular
Euclidean densities. The proof above is included because its densities
have compact support and zeros, and because the ambient sphere and the
finite-mixture path, rather than Euclidean regularity, are used here.

\begin{lemma}[Spectral entropy loss on a thin support]\label{lem:entropydefect61}
Suppose \eqref{eq:actiongap61} holds. Let $f$ be a bounded probability
density supported in a measurable set $A\subseteq\Sigma$ for which
\[
                         \Pi\mathbf1_A\le\theta<1.
\]
Then
\begin{equation}\label{eq:entropydefect61}
 \mathcal H(f)-\mathcal H(Pf)\ge(1-\kappa^2)(1-\theta).
\end{equation}
\end{lemma}
\begin{proof}
Orbital Cauchy--Schwarz gives
$(\Pi\sqrt f)^2\le(\Pi\mathbf1_A)(\Pi f)$ and hence
$\norm{\Pi\sqrt f}_2^2\le\theta$. The invariant and orthogonal
components are preserved by $P$. Since $\norm{\sqrt f}_2=1$,
\[
 \norm{P\sqrt f}_2^2
 \le\norm{\Pi\sqrt f}_2^2+\kappa^2\norm{(I-\Pi)\sqrt f}_2^2
 \le1-(1-\kappa^2)(1-\theta).
\]
For nonnegative numbers $z_a$ with $z=\sum_a\lambda_a z_a$, the
relative-entropy/Hellinger inequality gives
\[
 \sum_a\lambda_a z_a\log z_a-z\log z
 \ge\sum_a\lambda_a(\sqrt{z_a}-\sqrt z)^2.
\]
For completeness this follows by summing
$t\log(t/s)-t+s\ge(\sqrt t-\sqrt s)^2$; continuity handles zeros.
Apply it pointwise to $z_a=f(R(u_a)^{-1}x)$ and integrate. Invariance
of $\sigma$ and Cauchy--Schwarz give
\[
 \mathcal H(f)-\mathcal H(Pf)
 \ge2\left(1-\int\sqrt{Pf}\,P\sqrt f\,d\sigma\right)
 \ge2(1-\norm{P\sqrt f}_2).
\]
The inequality $2(1-\sqrt{1-d})\ge d$ for $0\le d\le1$ completes the
proof. No logarithmic Sobolev inequality or pointwise mixing estimate
has been assumed.
\end{proof}

\subsection{Transport paid for by loss of centroid mass}
For a unit physical random vector $X$ and a finite register $S$, put
\begin{equation}\label{eq:directionlaw61}
 c_s=\E[X\mid S=s],\quad
 \alpha=\sum_s\Pr(S=s)\norm{c_s},\quad
 \nu=\frac1\alpha\sum_{c_s\ne0}\Pr(S=s)\norm{c_s}\,
                       \delta_{c_s/\norm{c_s}}.
\end{equation}
Only positive-mass labels contribute. The directional law is defined
when $\alpha>0$; its weights are not, in general, the probabilities of
the hidden labels.

\begin{lemma}[Centroid-loss transport]\label{lem:centroidtransport61}
Apply an independent letter $a\sim\lambda$, followed by any fixed
physical identity-product word, using the actual stochastic rows of
the machine. Let $\alpha',\nu'$ be the quantities at the endpoint and
suppose $\alpha'\ge\zeta>0$. If $\lambda*\nu$ denotes the mixture of
the rotated laws $R(u_a)_*\nu$, then $\delta=\alpha-\alpha'\ge0$ and
\begin{equation}\label{eq:centroidtransport61}
 W_2(\lambda*\nu,\nu')^2\le\frac{3\pi^2}{\zeta}\delta.
\end{equation}
There is no bound on any intervening register.
\end{lemma}
\begin{proof}
Write $T_a(s,i)$ for the actual composite row, including the fixed
filler. Conditional independence of the row from the unrecorded past
is precisely the machine definition. Thus
\[
 p'_i c'_i=\sum_{s,a}p_s\lambda_aT_a(s,i)R(u_a)c_s.
\]
Set $w_{sai}=p_s\lambda_aT_a(s,i)\norm{c_s}$,
$q_i=\sum_{s,a}w_{sai}$, and $m_i=p'_i\norm{c'_i}$.
For $m_i>0$ put $v_i=c'_i/\norm{c'_i}$; otherwise choose an arbitrary
unit vector. With $u_{sa}=R(u_a)c_s/\norm{c_s}$ for nonzero terms,
\begin{equation}\label{eq:flowidentity61}
 \sum_{s,a}w_{sai}\norm{u_{sa}-v_i}^2=2(q_i-m_i).
\end{equation}
In particular $m_i\le q_i$, $\sum_iq_i=\alpha$,
$\sum_im_i=\alpha'$, and $\delta\ge0$.

Normalize the flows by $\alpha$. They couple $\lambda*\nu$ to
$\widehat\nu=\sum_i(q_i/\alpha)\delta_{v_i}$. Chordal and geodesic
distance satisfy $d_\Sigma(u,v)\le(\pi/2)\norm{u-v}$, so this
coupling has geodesic quadratic cost at most $\pi^2\delta/(2\alpha)$.
If $\delta>0$, the excess masses give a probability $\chi$ with
\[
 \widehat\nu=(\alpha'/\alpha)\nu'+(\delta/\alpha)\chi.
\]
Keeping the common part and transporting the remaining part across the
sphere costs at most $\pi^2\delta/\alpha$. The triangle inequality
for $W_2$ and $(a+b)^2\le2a^2+2b^2$ give
$W_2(\lambda*\nu,\nu')^2\le3\pi^2\delta/\alpha$.
Since $\alpha\ge\alpha'\ge\zeta$, this proves the result.
The case $\delta=0$ follows from the same couplings without an excess
measure. Zero centroids and zero-probability labels require no division.
\end{proof}

\subsection{The sharp lower law}
Assume, for some $A\ge1$, $p>0$, and $h_0>0$, the all-direction cap
bound
\begin{equation}\label{eq:allcaps61}
 \sup_{u,v\in\Sigma}m_G\{g:\norm{R(g)u-v}<h\}
                          \le Ah^p\quad(0<h\le h_0).
\end{equation}
This positive-power cap bound excludes every fixed unit vector. Hence
$E^G=\{0\}$ and the Haar mean of each orbit is zero; in particular
$0\in\operatorname{conv}R(G)z_0$. The constants are fixed with the
representation. For $k\ge1$ define
\begin{equation}\label{eq:parameters61}
 c_0=\min\{1,h_0,(2A)^{-1/p}\},\quad h_k=c_0k^{-1/p},\quad
 \gamma=(1-\kappa^2)/2,\quad L_k^{\rm ent}=\log(D_mh_k^{-(m-1)}).
\end{equation}

\begin{theorem}[Occupation without logarithmic loss]\label{thm:sharpoccupation61}
Let the command alphabet contain a physical identity letter, and assume
\eqref{eq:actiongap61} and \eqref{eq:allcaps61}. Fix a unit seed $z_0$,
let $Q=\operatorname{conv}R(G)z_0$, and suppose the target is
$\rho R(u_w)z_0$, with $0<\rho\le1$ and legal decoder values in $Q$.
Error is Euclidean error in $E$. For $0\le\varepsilon<\rho$ write
$\zeta=\rho-\varepsilon$. Every horizon-$N$ realization, including
one redesigned at that horizon, satisfies
\begin{equation}\label{eq:sharpoccupation61}
 \#\{1\le t\le N:|S_t|\le k\}\le
 \left\lceil 1+\frac{2L_k^{\rm ent}}\gamma+
       \frac{3\pi^2B_m}{\gamma^2h_k^2}\frac{1-\zeta}{\zeta}\right\rceil.
\end{equation}
In particular, for fixed subcritical parameters there are constants
$C_\zeta,c_\zeta>0$ such that
\begin{equation}\label{eq:sharpwidth61}
 \#\{1\le t\le N:|S_t|\le k\}\le C_\zeta k^{2/p},
                 \qquad W_{N,\varepsilon}\ge c_\zeta N^{p/2}.
\end{equation}
The intervening widths are unrestricted. The constants depend only on
$m,A,p,h_0,\kappa,\zeta$ and the fixed round normalization.
\end{theorem}
\begin{proof}
If there are $J\ge2$ positive cuts of width at most $k$, list them as
$t_1<\cdots<t_J$. Use a fixed prefix up to $t_1$. In each subsequent
gap put one independent letter with law $\lambda$, then identity
letters until the next selected cut. Finish by a fixed suffix. This is
an external distribution on actual words of length $N$, chosen from the
advertised width profile, not from the machine's private state. It does
not require the machine's rows on an identity letter to be identities.

Let $X$ be the current unit physical vector. Write $\alpha_j,\nu_j$
for \eqref{eq:directionlaw61} at $t_j$. Wordwise correctness and
$\norm{D_i}\le1$ imply at the terminal cut
\[
 \zeta\le\E\langle X_N,D(S_N)\rangle
       =\sum_i\Pr(S_N=i)\langle\E[X_N\mid S_N=i],D_i\rangle
       \le\alpha_N.
\]
The first inequality averages the conditional \emph{mean} error on
each deterministic word. It makes no assertion about a single sampled
label's error. Conditional Jensen, or the flow calculation above,
shows that centroid mass cannot increase along any gap or fixed
suffix. Thus $1\ge\alpha_j\ge\zeta$, and
\[
 \delta_j:=\alpha_j-\alpha_{j+1}\ge0,\qquad
                         \sum_{j<J}\delta_j\le1-\zeta.
\]

Fix $h=h_k$ and put $f_j=K_h\nu_j$. Each $f_j$ is supported on at most
$k$ open chordal caps of radius $h$. The cap assumption gives
$\Pi\mathbf1_{\{f_j>0\}}\le Akh^p\le1/2$.
Lemma~\ref{lem:entropydefect61} therefore gives
$\mathcal H(Pf_j)\le\mathcal H(f_j)-\gamma$.
Equivariance of the kernel says $Pf_j=K_h(\lambda*\nu_j)$.
Lemmas~\ref{lem:smoothentropy61} and~\ref{lem:centroidtransport61} imply
\[
 \mathcal H(f_{j+1})\le\mathcal H(f_j)-\gamma
              +\frac{\pi\sqrt{3B_m}}{h\sqrt\zeta}\sqrt{\delta_j}.
\]
Writing $M=J-1$, summing, using $0\le\mathcal H(f_j)\le L_h$, and
applying Cauchy--Schwarz only over these $M$ gaps, gives
\begin{equation}\label{eq:entropybudget61}
 \gamma M\le L_h+
       \frac{\pi\sqrt{3B_m}}h\sqrt{M\frac{1-\zeta}{\zeta}}.
\end{equation}
Young's inequality bounds the second term by $\gamma M/2+
3\pi^2B_m(1-\zeta)/(2\gamma h^2\zeta)$. This proves
\eqref{eq:sharpoccupation61}; $J\le1$ is immediate. The logarithm
in $L_h=O(1+\log k)$ is additive, not multiplied by $k^{2/p}$.
Since $1+\log k=O_p(k^{2/p})$, the first bound in
\eqref{eq:sharpwidth61} follows. If the whole profile has width at
most $k$, its $N$ positive cuts are all counted, proving the second.
\end{proof}

\begin{remark}[Endpoint information and synchronization]\label{rem:endbudget61}
The unrelaxed inequality \eqref{eq:entropybudget61} also holds.
When $\zeta=1$, it gives $J\le1+L_k^{\rm ent}/\gamma$, consistent
with exponential exact width for extreme outputs. For fixed positive
error it gives the polynomial law, but does not classify every joint
limit $N\to\infty$, $\varepsilon\downarrow0$.

If the identity letter is absent but physical identity-product words
exist at every length at least $g$, greedily retain narrow cuts at
spacing at least $g+1$. Each retained cut accounts for at most $g+1$
original narrow cuts. One independent command followed by a return of
the remaining length fits each retained gap. The prefix and suffix
may be any fixed words of their required lengths. The same proof
multiplies the occupation bound by at most $g+1$. Approximate returns
with accuracy-dependent thresholds are not included in this quantitative
extension without additional estimates.
\end{remark}

\begin{corollary}[Matching power laws]\label{cor:matchedorbits61}
In the setting of Theorem~\ref{thm:uniformorbit57}, replace its converse
by \eqref{eq:sharpwidth61}. In particular, if the all-direction exponent
$p$ equals the real dimension $q$ of the target orbit, then
\[
 W_{N,\varepsilon}=\Theta((N+1)^{q/2})
\]
for every fixed $0<\rho<1$, $0\le\varepsilon<\rho$, and also for
$\rho=1$, $0<\varepsilon<1$. For a general fixed-point-free compact
connected Lie representation one may take $p=p_*$ from
Theorem~\ref{thm:minorbit58}; the lower and upper exponents are then
$p_*/2$ and $q/2$, respectively, with no logarithmic loss in the lower
bound.
\end{corollary}
\begin{proof}
Apply Theorem~\ref{thm:sharpoccupation61}. The upper common-row
construction in Theorem~\ref{thm:uniformorbit57} is exact for
$\rho<1$ and has $O((N+1)^{q/2})$ labels. At $\rho=1$ and positive
error use its amplitude-slack version. The two bounds match when
$p=q$. Theorem~\ref{thm:minorbit58} supplies the remaining assertion.
\end{proof}
The endpoint $\rho=1,\varepsilon=0$ is not an assertion of the matching
polynomial corollary. No exact leading constant or asymptotic ratio is
claimed. The spectral constant is effective only when a certified value
is supplied; a fully numerical example is given below.
